\chapter{Structure du marché américain des actions}\label{ch:m1:us-equity-market-structure}

Cent actions d'une grande société américaine peuvent être achetées à cet instant
sur plus d'une douzaine de bourses, dans des dizaines de \omterm{def:m1:us-equity-market-structure:dark}{dark pools} et
auprès de dealers qui ne montrent jamais de cotation à personne. Plus de la moitié des
17.6 milliards d'actions échangées lors d'une séance moyenne de 2025 ne l'ont été
sur aucune des bourses. Pendant deux décennies, une seule réglementation a relié ces
lieux entre eux en décrétant lequel de leurs prix est «~le~» prix et en interdisant
les transactions à de moins bons prix~; en juin 2026, le régulateur a proposé d'abroger
ce décret. Ce chapitre décrit la machine telle qu'elle existe aujourd'hui~: les
règles, la cotation consolidée qu'elles créent, les deux vitesses auxquelles cette
cotation peut être connue, les jeux de frais qui l'entourent, et la moitié du
marché qui vit en dehors des bourses.

\section{Regulation NMS}

La Securities and Exchange Commission a adopté la Regulation National Market
System en 2005. Quatre de ses règles régissent chaque ordre.

\begin{definition}[Cotation protégée, trade-through, règle de protection des ordres]\label{def:m1:us-equity-market-structure:opr}
Une \emph{cotation protégée} (protected quotation) est le meilleur cours acheteur ou le
meilleur cours vendeur d'une bourse, affiché, automatisé et immédiatement accessible. Une
\emph{trade-through}\index{trade-through} est l'exécution d'une transaction à un
prix inférieur à une cotation protégée affichée ailleurs~: un achat au-dessus de
le cours vendeur protégé d'une autre plateforme, ou une vente sous son cours acheteur protégé. La
\emph{règle de protection des ordres}\index{règle de protection des ordres} (order protection rule, Rule 611) impose à
chaque centre de négociation de disposer de politiques raisonnablement conçues pour
empêcher les trade-throughs.
\end{definition}

\begin{definition}[Plafond des frais d'accès]\label{def:m1:us-equity-market-structure:cap}
Le \emph{plafond des frais d'accès}\index{plafond des frais d'accès} (access fee cap, Rule 610(c)) limite ce qu'une
plateforme peut facturer pour exécuter contre sa cotation protégée~: 0.003~\$ par
action («~30 mils~») pour les titres valant un dollar ou plus. Sans elle, une
plateforme pourrait afficher un prix attractif que la règle de protection oblige les autres
à router vers elle, et facturer ce qu'elle voulait pour l'atteindre.
\end{definition}

La Rule 610 interdit aussi les cotations qui \emph{verrouillent} ou \emph{croisent} le marché
(un cours acheteur égal ou supérieur au cours vendeur d'une autre plateforme). La Rule 612, la
\emph{règle des sous-cents} (sub-penny rule), fixe le pas de cotation minimal à un cent pour
les titres valant un dollar ou plus. La Rule 603 régit la diffusion des données de marché,
auxquelles nous revenons plus bas.

\begin{dated}{2026-09}{Regulation NMS en mouvement}\label{dat:m1:us-equity-market-structure:nms}
En septembre 2024, la SEC a adopté des amendements créant un pas de cotation d'un demi-cent
pour les titres «~contraints par le tick~», et ramenant le \omterm{def:m1:us-equity-market-structure:cap}{plafond des frais d'accès} à
0.001~\$. Aucun des deux n'est en vigueur~: par une ordonnance du 11 juin 2026, la mise en
application de ces deux dispositions est reportée au premier jour ouvré de novembre 2027. Le même jour,
la SEC a \emph{proposé d'abroger} la Rule 611 ainsi que l'interdiction des marchés
verrouillés et croisés, dans leur totalité, avec une période de commentaires de 60 jours.
Le reste de ce chapitre décrit les règles en vigueur~; il faut consulter l'état actuel de la
proposition avant de construire quoi que ce soit qui en dépende.
\end{dated}

\section{Les meilleurs cours acheteur et vendeur nationaux}

\begin{definition}[Lot standard et lot irrégulier]\label{def:m1:us-equity-market-structure:lots}
Un \emph{lot standard}\index{lot standard} (round lot) est l'unité de négociation standard qu'une
cotation doit atteindre pour être protégée~; un \emph{lot irrégulier}\index{lot irrégulier} (odd lot) est
toute quantité plus petite.
\end{definition}

\begin{dated}{2026-09}{Lots standard selon le prix}\label{dat:m1:us-equity-market-structure:lots}
Depuis le 3 novembre 2025, le \omterm{def:m1:us-equity-market-structure:lots}{lot standard} dépend du cours de clôture moyen de l'action
sur la période d'évaluation précédente~: 100 actions jusqu'à 250.00~\$~; 40 actions de
250.01~\$ à 1\,000.00~\$~; 10 actions de 1\,000.01~\$ à
10\,000.00~\$~; 1 action au-delà. À partir de la même date, les flux consolidés
indiquent les tailles en actions et non plus en lots. Jusqu'alors, un cours acheteur de 50 actions
d'un titre à 3\,000~\$, valant 150\,000~\$, était un \omterm{def:m1:us-equity-market-structure:lots}{lot irrégulier}~: il n'était ni protégé ni
inclus dans le NBBO.
\end{dated}

\begin{definition}[Meilleurs cours acheteur et vendeur nationaux]\label{def:m1:us-equity-market-structure:nbbo}
Les \emph{meilleurs cours acheteur et vendeur nationaux}\index{meilleurs cours acheteur et vendeur nationaux}
(national best bid and offer, NBBO) d'un titre sont le cours acheteur protégé le plus élevé et le cours vendeur protégé le
plus bas, sur l'ensemble des bourses, avec la taille totale affichée à chacun de ces
prix.
\end{definition}

\begin{example}[Trois plateformes]\label{ex:m1:us-equity-market-structure:nbbo}
La plateforme X affiche $10.00 \times 10.02$ (300 par 500 actions), la plateforme Y $10.01
\times 10.02$ (200 par 100), la plateforme Z $10.01 \times 10.03$ (400 par 900). Le
NBBO est $10.01 \times 10.02$, 600 par 600~: le cours acheteur de Y et de Z, le cours vendeur
de X et de Y. Un courtier qui exécute l'achat au marché d'un client à 10.03 alors que
X et Y offrent à 10.02 a provoqué un \omterm{def:m1:us-equity-market-structure:opr}{trade-through} sur deux cotations protégées. Si une
quatrième plateforme affiche un cours acheteur de 10.015 pour 60 actions, rien ne
change~: un \omterm{def:m1:us-equity-market-structure:lots}{lot irrégulier} n'est pas protégé. Le NBBO est l'indice de référence de toutes les
statistiques de qualité d'exécution du \cref{ch:m1:retail-flow-and-wholesaling}~; il n'est
pas non plus une seule chose, comme nous le voyons maintenant.
\end{example}

\section{Un marché, deux horloges}

\begin{definition}[Processeur d'information sur les titres et flux direct]\label{def:m1:us-equity-market-structure:sip}
Un \emph{processeur d'information sur les titres}\index{processeur d'information sur les titres}
(securities information processor, SIP) collecte les meilleures cotations et les transactions de chaque bourse, calcule le
NBBO et publie le flux consolidé. Un
\emph{flux direct}\index{flux direct} (direct feed) est le flux de données propriétaire vendu par
une bourse, contenant l'intégralité de son carnet d'ordres, à partir duquel un abonné peut
calculer sa propre vue consolidée.
\end{definition}

\begin{omfigure}
\begin{tikzpicture}[font=\small,
  ven/.style={draw=omExo,rounded corners=2pt,minimum width=2.1cm,minimum height=0.7cm,fill=omExo!6},
  box/.style={draw=omDef,rounded corners=2pt,minimum width=3.2cm,minimum height=0.9cm,align=center,fill=omDef!4},
  slow/.style={-{Stealth[length=2mm]},thick,omThm},
  fast/.style={-{Stealth[length=2mm]},thick,omDef}]
\foreach \i/\y/\n in {a/1.8/Bourse A,b/0.6/Bourse B,c/-0.6/Bourse C,d/-1.8/Bourse D}{ \node[ven] (v\i) at (0,\y) {\n}; }
\node[box] (sip) at (5.6,1.5) {SIP\\(consolide, republie)};
\node[box] (own) at (5.6,-1.5) {Constructeur de carnet propre de la société\\dans les mêmes centres de données};
\node[box] (slowu) at (11.6,1.5) {Abonné au flux consolidé\\voit le NBBO};
\node[box] (fastu) at (11.6,-1.5) {Abonné au flux direct\\le voit plus tôt};
\foreach \i in {a,b,c,d}{ \draw[slow] (v\i.east) -- (sip.west); \draw[fast] (v\i.east) -- (own.west); }
\draw[slow] (sip) -- node[above,font=\footnotesize]{$+\,\delta$} (slowu);
\draw[fast] (own) -- (fastu);
\end{tikzpicture}

\omcaption{Les mêmes cotations, consolidées deux fois. Le parcours de la bande (en rouge) ajoute
un retard $\delta$~: le trajet jusqu'au processeur, son calcul et le trajet
jusqu'à l'abonné. Le chemin direct (en bleu) est aussi rapide que les
lignes et le code propres à l'abonné.}\label{fig:m1:us-equity-market-structure:twoclocks}
\end{omfigure}

\begin{proposition}[À quelle fréquence la bande est-elle fausse]\label{prop:m1:us-equity-market-structure:stale}
Supposons que le meilleur cours acheteur change aux instants d'un processus de Poisson d'intensité
$\lambda$, et la bande rapporte chaque changement avec un retard $\delta$ par rapport au moment où un lecteur du \omterm{def:m1:us-equity-market-structure:sip}{flux direct}
l'observe. Si $\lambda\delta \ll 1$, la fraction du temps pendant laquelle
les deux vues diffèrent vaut approximativement $\lambda\delta$.
\end{proposition}
\begin{proof}
Chaque changement ouvre un intervalle de longueur $\delta$ pendant lequel la bande
affiche encore l'ancien cours acheteur. Lorsque $\lambda\delta \ll 1$, ces intervalles
se chevauchent avec une probabilité négligeable, et sur une période $T$ ils couvrent $\lambda T \cdot
\delta$ en espérance.
\end{proof}

\begin{example}[La moitié d'un pour cent du temps, toutes les transactions qui comptent]\label{ex:m1:us-equity-market-structure:halfpercent}
Avec cinq changements par seconde et $\delta = 1$ milliseconde, la bande est périmée
0.5~\% du temps~; lors d'une rafale de cinquante changements par seconde, 5~\%
(\cref{fig:m1:us-equity-market-structure:stale}). La fraction paraît négligeable
et ne l'est pas~: ce sont précisément les moments où le prix évolue, et un ordre fixé
d'après la vue périmée est exécuté \emph{parce que} celle-ci est
périmée. Une cotation qui n'a pas encore suivi un mouvement de deux centimes peut être achetée un
centime en dessous de la nouvelle fair value~; dans l'action simulée du tutoriel, les
plateformes plus lentes offrent 3\,636 occasions de ce type en dix minutes. Qui a le droit de les saisir
et comment les plateformes protègent leurs participants lents, relève de One
Quant Book 11.
\end{example}

\begin{omfigure}
\begin{tikzpicture}
\begin{axis}[width=10.5cm,height=4.8cm,
  xlabel={délai de la bande $\delta$ (microsecondes)},ylabel={bid consolidé périmé (\% du temps)},
  tick label style={font=\small},label style={font=\small},
  xmin=0,xmax=2000,ymin=0,ymax=10,grid=major,grid style={gray!25},xtick={0,500,1000,1500,2000},
  xticklabel style={/pgf/number format/1000 sep={\,}},
  legend columns=2,legend style={font=\footnotesize,at={(0.5,-0.34)},anchor=north,draw=none,fill=none,/tikz/every even column/.append style={column sep=8pt}},
  table/col sep=comma]
\addplot[omThm,very thick,mark=*,mark size=1.2pt] table[x=sip_delay_us,y=stale_pct_busy]{figdata/markets-1/09-us-equity-market-structure/stale.csv};
\addplot[omDef,very thick,mark=*,mark size=1.2pt] table[x=sip_delay_us,y=stale_pct_quiet]{figdata/markets-1/09-us-equity-market-structure/stale.csv};
\legend{50 changements de prix par seconde,5 changements de prix par seconde}
\end{axis}
\end{tikzpicture}

\omcaption{Péremption du meilleur cours acheteur consolidé en fonction du délai de la bande, pour
quatre plateformes qui recotent 50 à 400 microsecondes après chaque mouvement. Les pentes sont
les deux valeurs de $\lambda$. Données~: la
simulation du tutoriel.}\label{fig:m1:us-equity-market-structure:stale}
\end{omfigure}

\section{Familles de bourses et jeu des frais}

La plupart des bourses appartiennent à quelques groupes, chacun exploitant plusieurs
plateformes avec des grilles tarifaires différentes sur une technologie presque identique. La raison
en est le \omterm{def:m1:us-equity-market-structure:cap}{plafond des frais d'accès}~: les plateformes ne peuvent pas rivaliser sur le prix affiché, que le
tick fixe à un cent, elles rivalisent donc à l'intérieur du plafond, sur qui paie quoi à qui.

\begin{definition}[Tarification maker-taker et plateforme inversée]\label{def:m1:us-equity-market-structure:makertaker}
Sous \emph{tarification maker-taker}\index{tarification maker-taker}, une plateforme facture
des frais à l'ordre qui retire de la liquidité (le \emph{preneur de liquidité}) et verse
un rabais plus faible à l'ordre au repos contre lequel il s'exécute (le \emph{fournisseur de liquidité}),
et garde la différence. Une \emph{plateforme inversée}\index{plateforme inversée} fait
l'inverse~: elle paie le preneur et facture le fournisseur.
\end{definition}

\begin{proposition}[Les frais équivalent à un tick plus fin]\label{prop:m1:us-equity-market-structure:finer}
Soient deux plateformes affichant le même cours vendeur $a$. Un acheteur qui prend de la liquidité paie
$a + \varphi$ sur une plateforme facturant des frais preneurs $\varphi$, et $a - \rho$ sur une
\omterm{def:m1:us-equity-market-structure:makertaker}{plateforme inversée} offrant un rabais preneur $\rho$. Un preneur rationnel s'adresse donc d'abord à la
\omterm{def:m1:us-equity-market-structure:makertaker}{plateforme inversée}~: un fournisseur qui y affiche un ordre est en tête
d'une file \emph{intermarchés}, moyennant un prix, et reçoit $a - \mu$, où $\mu$ désigne
les frais fournisseur de la \omterm{def:m1:us-equity-market-structure:makertaker}{plateforme inversée}.
\end{proposition}
\begin{proof}
Immédiat à partir des définitions~; l'ordre des plateformes selon leur prix tout compris
$a + \text{frais}$ est strict bien que les prix affichés soient égaux.
\end{proof}

\begin{example}[Priorité à l'achat]\label{ex:m1:us-equity-market-structure:priority}
Avec des frais illustratifs dans le plafond~: sur une plateforme maker-taker, un vendeur
au repos à 10.00 perçoit un rabais de 0.20 cent mais attend derrière tous ceux qui
sont arrivés plus tôt~; sur une \omterm{def:m1:us-equity-market-structure:makertaker}{plateforme inversée}, il paie 0.18 cent et est exécuté avant
toutes les plateformes maker-taker, car les preneurs touchent 0.15 cent pour s'y rendre
en premier. La différence, 0.38 cent, est le prix de la priorité sur un
marché dont le tick vaut un cent. Un \omterm{def:m1:what-a-trading-firm-does:market-maker}{teneur de marché} choisit par action et par
moment~: là où la file est longue et où la fourchette vaut un tick, la priorité vaut
la peine d'être achetée (\cref{fig:m1:us-equity-market-structure:fees}).
\end{example}

\begin{omfigure}
\begin{tikzpicture}
\begin{axis}[width=9.2cm,height=3.8cm,xbar,bar width=8pt,
  xlabel={coût tout compris d'une exécution au cours vendeur, par rapport à la cotation (cents)},
  ytick=data,yticklabels={{Maker-taker (frais de prise 0.30 c)},{Frais fixes (0.09 c)},{Inversée (rebate de prise 0.15 c)}},
  yticklabel style={font=\small},tick label style={font=\small},label style={font=\small},
  y dir=reverse,xmin=-0.35,xmax=0.45,xtick={-0.3,-0.15,0,0.15,0.3,0.45},enlarge y limits=0.25,grid=major,grid style={gray!25},
  xticklabel style={/pgf/number format/fixed,/pgf/number format/precision=2},
  nodes near coords,nodes near coords style={font=\footnotesize,/pgf/number format/fixed,/pgf/number format/precision=2},
  table/col sep=comma]
\addplot[fill=omDef!60,draw=omDef] table[x=net_cents_vs_quote,y=k]{figdata/markets-1/09-us-equity-market-structure/fees.csv};
\end{axis}
\end{tikzpicture}

\omcaption{Trois plateformes affichant le même cours vendeur sont trois prix différents
pour un preneur de liquidité. Les niveaux de frais sont illustratifs, dans le plafond de 0.30 cent. Données~: le
script du chapitre.}\label{fig:m1:us-equity-market-structure:fees}
\end{omfigure}

\begin{definition}[Ordre de balayage intermarchés]\label{def:m1:us-equity-market-structure:iso}
Un \emph{ordre de balayage intermarchés}\index{ordre de balayage intermarchés} (ISO) est un
ordre à cours limité marqué de sorte que la plateforme réceptrice puisse l'exécuter sur-le-champ sans
vérifier les cotations protégées des autres plateformes, parce que l'expéditeur a
simultanément acheminé des ordres pour consommer toutes les cotations protégées
de meilleur prix ailleurs. La conformité à la \omterm{def:m1:us-equity-market-structure:opr}{règle de protection des ordres} passe
de la plateforme à l'expéditeur.
\end{definition}

L'ISO est la façon dont tout intervenant de grande taille ou rapide négocie en réalité~: il envoie, dans
la même microseconde, un ordre par plateforme, dimensionné sur la quantité affichée,
et balaie ainsi plusieurs niveaux de prix sur le marché avant qu'aucune plateforme ne puisse
réacheminer quoi que ce soit. Sans lui, un ordre dépassant la taille offerte au meilleur
prix devrait traverser les plateformes au rythme d'une décision de routage à la fois.

\section{La part hors bourse}

\begin{definition}[Dark pool]\label{def:m1:us-equity-market-structure:dark}
Un \emph{dark pool}\index{dark pool} est un \omterm{def:m1:exchanges-brokers-venues:ats}{système de négociation alternatif}
(la \cref{def:m1:exchanges-brokers-venues:ats}) qui n'affiche aucune cotation~:
les ordres restent invisibles et s'exécutent les uns contre les autres à des prix dérivés du
NBBO, le plus souvent son point médian.
\end{definition}

Les transactions réalisées hors bourse restent publiques~: elles sont déclarées dans
quelques secondes à un \emph{dispositif de déclaration des transactions} et figurent dans le flux des transactions sans
le nom du lieu. Elles proviennent de deux sources. Les \omterm{def:m1:us-equity-market-structure:dark}{dark pools} servent
des institutions qui ne veulent pas dévoiler un ordre de grande taille et se contentent
d'attendre une contrepartie au point médian. Les \emph{dealers en principal} --- wholesalers
exécutant les ordres des courtiers de détail contre leur propre inventaire
(du \cref{ch:m1:retail-flow-and-wholesaling}) et les plateformes à dealer unique des banques
--- interviennent en tant que contrepartie, non en tant que plateforme.

\begin{dated}{2026-09}{Où se sont échangées les actions américaines en 2025}\label{dat:m1:us-equity-market-structure:where}
Le volume quotidien moyen en 2025 s'est élevé à 17.6~milliards d'actions, pour une valeur de 1.1~billion de dollars, en hausse
d'environ 45~\% par rapport à 2024. La part des dispositifs de déclaration des transactions dans le volume consolidé
a atteint 50.6~\%, au-dessus de la moitié pour la première fois sur une année pleine~;
81.3~\% de ce volume hors bourse ont été exécutés par des dealers en principal et
18.7~\% sur des \omterm{def:m1:exchanges-brokers-venues:ats}{systèmes de négociation alternatifs}, selon le rapport annuel
d'un groupe boursier.
\end{dated}

\begin{omfigure}
\begin{tikzpicture}
\begin{axis}[width=9cm,height=3.8cm,xbar,bar width=8pt,
  xlabel={part du volume des actions américaines en 2025 (\%)},
  ytick=data,yticklabels={{Bourses},{Hors bourse~: dealers en principal},{Hors bourse~: systèmes de négociation alternatifs}},
  yticklabel style={font=\small},tick label style={font=\small},label style={font=\small},
  y dir=reverse,xmin=0,xmax=60,enlarge y limits=0.25,grid=major,grid style={gray!25},
  nodes near coords,nodes near coords style={font=\footnotesize},
  table/col sep=comma]
\addplot[fill=omDef!60,draw=omDef] table[x=pct,y=k]{figdata/markets-1/09-us-equity-market-structure/where_2025.csv};
\end{axis}
\end{tikzpicture}

\omcaption{Où une action a été échangée en 2025. Les bourses, qui produisent les
cotations qui fixent le prix de tout le reste, n'ont exécuté qu'une minorité du volume. Données~:
comme dans l'\cref{dat:m1:us-equity-market-structure:where}.}
\end{omfigure}

\begin{remark}[Pourquoi cette répartition inquiète les régulateurs]\label{rem:m1:us-equity-market-structure:free}
Les transactions hors bourse sont fixées à partir du NBBO et n'y contribuent en rien~:
elles profitent gratuitement des cotations affichées. Si les ordres les moins dangereux pour un
\omterm{def:m1:what-a-trading-firm-does:market-maker}{teneur de marché} --- les ordres de détail --- sont exécutés hors des bourses, les ordres
qui atteignent les bourses sont en moyenne mieux informés, et les fourchettes des bourses
s'élargissent (la \cref{prop:m1:what-a-trading-firm-does:capture}), et l'indice de référence
par rapport auquel on mesure l'«~amélioration de prix~» hors bourse devient plus facile
à battre. Savoir si cette boucle est importante en pratique est l'une des questions
empiriques les plus disputées de la structure de marché.
\end{remark}

\section{Ralentisseurs}

Si la \cref{prop:m1:us-equity-market-structure:stale} décrit un préjudice, une plateforme
peut vendre une protection contre celui-ci. En 2016, la SEC a approuvé en tant que bourse une plateforme,
IEX, dont chaque message entrant et sortant transite par 38~miles de
fibre optique enroulée, soit un délai de 350~microsecondes, tandis que la vision propre
de la plateforme sur le marché n'est \emph{pas} retardée. Les ordres au repos indexés sur le point médian
sont donc recotés par la plateforme avant qu'un trader rapide ayant vu le
même mouvement puisse les atteindre. La SEC a jugé qu'un délai d'une telle ampleur est
\emph{de minimis}~: les cotations de la plateforme restent «~immédiatement accessibles~»
et protégées. Le débat sur la question de savoir qui doit supporter le coût de la vitesse est loin
d'être clos.

\section{Tutoriel~: construire un NBBO et mesurer le retard du flux}\label{tut:m1:us-equity-market-structure:nbbo}

\begin{tutorial}
\textbf{Objectif.} Consolider les cotations de chaque plateforme en un NBBO, signaler les marchés verrouillés et
croisés et les \omterm{def:m1:us-equity-market-structure:opr}{trade-through}, puis mesurer à quel point un flux retardé est périmé.
\textbf{État final~:} les résultats de
l'\cref{ex:m1:us-equity-market-structure:nbbo} et
la \cref{fig:m1:us-equity-market-structure:stale}.
\begin{tutsteps}
\item \textbf{Le constructeur}, tel qu'il s'exécute dans le processus de trading de la société~: une
table fixe des plateformes, sans allocation, une seule passe par mise à jour.
\omcode{code/firm/nbbo/cpp/firm_nbbo.hpp}{41}{58}{Recalcul du NBBO à partir de la table des plateformes~: seules les tailles protégées, égalités regroupées.}\label{lst:m1:us-equity-market-structure:cpp}
\item \textbf{Le vérifier sur trois plateformes.} Les tests d'acceptation rejouent
l'\cref{ex:m1:us-equity-market-structure:nbbo}~: meilleur cours acheteur de 10.01 pour 600 actions
sur deux plateformes~; un cours acheteur de 10.015 en \omterm{def:m1:us-equity-market-structure:lots}{lot irrégulier}, ignoré jusqu'à ce qu'il atteigne un \omterm{def:m1:us-equity-market-structure:lots}{lot standard}~;
un achat à 10.03 signalé comme \omterm{def:m1:us-equity-market-structure:opr}{trade-through}.
\item \textbf{Simuler deux horloges.} Des sauts de un à trois ticks arrivent à
instants de Poisson~; chaque plateforme recote après son propre délai~; le flux ajoute
$\delta$.
\omcode{code/markets-1/09-us-equity-market-structure/python/stale_tape.py}{29}{49}{Cours de chaque plateforme en fonction du temps, sous forme de fonction en escalier, et mesure de l'ensemble où le flux et la vue directe diffèrent.}\label{lst:m1:us-equity-market-structure:stale}
\item \textbf{Comparer avec la proposition.} Pour $\delta = 1\,000$
microsecondes et $\lambda = 5$ par seconde, la simulation donne 0.51~\%
contre 0.50~\% prévus~; pour $\lambda = 50$, 4.9~\% contre 5~\%, l'écart
étant dû aux chevauchements que la proposition néglige.
\end{tutsteps}
\textbf{Pour aller plus loin.} Donner à une plateforme un ralentisseur de 350~microsecondes sur
les seuls ordres entrants, et compter combien de ses cotations périmées survivent. Puis
ajouter une \omterm{def:m1:us-equity-market-structure:makertaker}{plateforme inversée} et acheminer un preneur de liquidité selon le prix tout compris.
\end{tutorial}

\section{Construction~: la cotation consolidée}\label{bld:m1:us-equity-market-structure:nbbo}

\begin{build}
\bfield{Rôle} Chaque stratégie, chaque routeur et chaque rapport d'exécution de la société miniature
a besoin des meilleurs cours acheteur et vendeur entre les plateformes, calculés par la société
elle-même à partir des \omterm{def:m1:us-equity-market-structure:sip}{flux directs} (les gestionnaires de flux sont construits dans le
livre~13).
\bfield{Interface} \texttt{NbboBuilder(round\_lot)}~;
\texttt{update(venue, bid, bid\_size, ask, ask\_size)} renvoyant le nouveau NBBO
s'il a changé~; \texttt{nbbo()}~; \texttt{trades\_through(side, price)}. Le
NBBO porte les prix, les tailles agrégées, l'ensemble des plateformes à chaque prix, et
les indicateurs \texttt{locked} et \texttt{crossed}.
\bfield{Règles} Prix entiers. Une cotation inférieure au \omterm{def:m1:us-equity-market-structure:lots}{lot standard} est stockée
mais non protégée. Un côté manquant signifie qu'il n'y a pas de NBBO. En C++ et en Rust~: une
table de plateformes fixe, aucune allocation ni verrou sur le chemin de mise à jour~; les plateformes
sont de petits entiers et leurs ensembles sont des masques de bits.
\bfield{Tests d'acceptation} \texttt{code/firm/nbbo/tests/}, \texttt{cpp/} et
\texttt{rust/}~: l'exemple à trois plateformes~; les \omterm{def:m1:us-equity-market-structure:lots}{lots irréguliers}~; la détection de changement~;
les cours verrouillés et croisés~; le \omterm{def:m1:us-equity-market-structure:opr}{trade-through}~; le rejet des cotations mal formées.
\bfield{Pour aller plus loin} \omterm{def:m1:us-equity-market-structure:lots}{Lots standard} par symbole, à partir d'une table de référence
(voir l'\cref{dat:m1:us-equity-market-structure:lots})~; un second NBBO, \emph{ajusté des frais},
calculé avec les frais preneur de chaque plateforme
(voir la \cref{prop:m1:us-equity-market-structure:finer}).
\end{build}

\begin{omsources}
\item US Securities and Exchange Commission, \emph{Regulation NMS}, Release
34-51808, \emph{Federal Register} 70 (29 juin 2005).
\item US Securities and Exchange Commission, \emph{Fact Sheet: Regulation NMS
Reforms} (proposition du 11 juin 2026, Release 34-105655)~; Release 34-105656,
ordonnance accordant une dispense temporaire, 11 juin 2026~; communiqué de presse
2024-137 (modifications du 18 septembre 2024).
\item Cboe Global Markets, \emph{Regulation NMS Round Lots Enhancements FAQ},
28 octobre 2025~; \emph{2025 U.S. Equities Year in Review}.
\item US Securities and Exchange Commission, Release 34-78102,
\emph{Commission Interpretation Regarding Automated Quotations Under
Regulation NMS}, juin 2016~; IEX, \emph{The speed bump}.
\item E.~Budish, P.~Cramton et J.~Shim, ``The high-frequency trading arms
race'', \emph{Quarterly Journal of Economics} 130 (2015).
\item M.~O'Hara, ``High frequency market microstructure'', \emph{Journal of
Financial Economics} 116 (2015).
\end{omsources}

\section{Exercices}

\begin{exercise}[$\star$]\label{exo:m1:us-equity-market-structure:1}
Quatre plateformes cotent un titre avec un \omterm{def:m1:us-equity-market-structure:lots}{lot standard} de 100 actions~: A $25.10 \times 25.14$
(500 par 200), B $25.11 \times 25.13$ (100 par 300), C $25.11 \times 25.13$ (300
par 50), D $25.12 \times 25.15$ (80 par 400). Donner le NBBO avec les tailles et
les plateformes.
\end{exercise}

\begin{exercise}[$\star$]\label{exo:m1:us-equity-market-structure:2}
En utilisant l'\cref{dat:m1:us-equity-market-structure:lots}, donner le \omterm{def:m1:us-equity-market-structure:lots}{lot standard} des
titres cotés 48~\$, 251~\$, 999~\$, 4\,200~\$ et 712\,000~\$, ainsi que la valeur en dollars d'un
\omterm{def:m1:us-equity-market-structure:lots}{lot standard} de chacun.
\end{exercise}

\begin{exercise}[$\star$]\label{exo:m1:us-equity-market-structure:3}
À partir de l'\cref{dat:m1:us-equity-market-structure:where}, calculer le nombre quotidien moyen
d'actions exécutées sur les bourses, par les dealers agissant en principal et sur les \omterm{def:m1:exchanges-brokers-venues:ats}{systèmes
de négociation alternatifs}, ainsi que le prix moyen d'une action échangée.
\end{exercise}

\begin{exercise}[$\star\star$]\label{exo:m1:us-equity-market-structure:4}
Avec le NBBO de l'exercice~1, un courtier exécute un achat de 200 actions pour
un client au prix de 25.14 sur la plateforme A. Quelles cotations protégées ont été traversées~?
Qu'aurait-il dû envoyer, et avec quelle mention, pour acheter 500 actions immédiatement et
en toute légalité~?
\end{exercise}

\begin{exercise}[$\star\star$]\label{exo:m1:us-equity-market-structure:5}
Une plateforme maker-taker facture 0.30 cent aux preneurs de liquidité et rémunère les \omterm{def:m1:us-equity-market-structure:makertaker}{fournisseurs
de liquidité} 0.25~; une \omterm{def:m1:us-equity-market-structure:makertaker}{plateforme inversée} rémunère les preneurs de 0.10 et facture 0.14 aux
fournisseurs. Les deux affichent un cours vendeur de 40.00. Donner le prix tout compris payé par un
preneur et reçu par un fournisseur sur chacune, le revenu par action de la plateforme sur chacune et le prix de la priorité.
\end{exercise}

\begin{exercise}[$\star\star$]\label{exo:m1:us-equity-market-structure:6}
Le meilleur cours acheteur d'un titre change 12 fois par seconde en moyenne. Estimer la fraction de
temps pendant laquelle un flux retardé de 600 microsecondes affiche un cours acheteur périmé, et
le nombre de secondes par journée de négociation de 6.5 heures. Pourquoi la proposition surestime-t-elle
légèrement lorsque les changements se regroupent~?
\end{exercise}

\begin{exercise}[$\star\star\star$]\label{exo:m1:us-equity-market-structure:7}
\emph{Programmation.} Avec \texttt{stale\_tape}, diviser par deux le délai de recotation de chaque
plateforme et recalculer la fraction périmée pour $\delta = 500$ microsecondes et $\lambda =
5$. Expliquer pourquoi le résultat change à peine, et quelle grandeur les délais des plateformes
\emph{contrôlent} effectivement.
\end{exercise}

\begin{exercise}[$\star\star\star$]\label{exo:m1:us-equity-market-structure:8}
\emph{Trouver le défaut.} Le rapport d'un courtier de détail indique~: «~98~\% des ordres de nos clients
ont été exécutés au NBBO ou à un meilleur prix.~» Le courtier calcule le
NBBO à partir du flux consolidé et horodate chaque ordre à sa réception sur ses propres serveurs.
Donner deux raisons pour lesquelles cette statistique surestime la qualité d'exécution, et une
raison propre aux titres à prix élevé avant novembre 2025.
\end{exercise}

\section{Problème~: un flux périmé}

\begin{problem}[{Problème du week-end --- deux vues d'un même titre}]\label{pb:m1:us-equity-market-structure:1}
Un titre est négocié sur quatre bourses. Chacune cote un tick (un centime) de part et d'autre
de son estimation du prix juste, pour 300 actions par côté. Lorsque le prix juste saute, les
\omterm{def:m1:what-a-trading-firm-does:market-maker}{teneurs de marché} des bourses recotent après 60, 150, 250 et 500
microsecondes respectivement. Le flux consolidé parvient à un abonné 900 microsecondes
après celui d'un \omterm{def:m1:us-equity-market-structure:sip}{flux direct}. Le prix juste saute 8 fois par minute en moyenne~;
65~\% des sauts sont d'un tick, 25~\% de deux ticks et 10~\% de trois.

\medskip\noindent\textbf{Partie I --- Le flux.}
\begin{enumerate}
\item Quelle fraction du temps le flux affiche-t-il un meilleur cours acheteur périmé~?
\item Combien de secondes cela représente-t-il sur une journée de 6.5 heures~?
\item Après un saut à la hausse, la mise à jour de quelle plateforme déplace le meilleur cours acheteur
sur un \omterm{def:m1:us-equity-market-structure:sip}{flux direct}, et au bout de combien de temps~?
\item Après le même saut, quand le \emph{cours vendeur} bouge-t-il, et pourquoi
la réponse est-elle différente~?
\end{enumerate}

\medskip\noindent\textbf{Partie II --- Le snipe.}
Un trader rapide voit chaque \omterm{def:m1:us-equity-market-structure:sip}{flux direct} et peut atteindre n'importe quelle plateforme
40 microsecondes après un saut.
\begin{enumerate}[resume]
\item Après un saut à la hausse de $J$ ticks, quelles plateformes affichent encore
l'ancien cours vendeur lorsque le trader arrive~?
\item En achetant à un ancien cours vendeur et en valorisant au nouveau prix juste, quel
est le profit par action pour $J = 1, 2, 3$~?
\item Donner le profit espéré par action, par saut et par plateforme périmée.
\item Avec 300 actions sur chaque plateforme périmée, donner le profit espéré par saut
et par jour si le trader gagne toutes les courses.
\item Le trader paie des frais de preneur de 0.30 centime. Recalculer le profit espéré
par action, par saut et par plateforme. Quels sauts valent-ils encore la peine d'être pris~?
\end{enumerate}

\medskip\noindent\textbf{Partie III --- La défense.}
\begin{enumerate}[resume]
\item Le \omterm{def:m1:what-a-trading-firm-does:market-maker}{teneur de marché} le plus lent perd ce que le trader gagne. Combien lui coûte
le fait d'être victime d'un snipe, par jour et par action échangée s'il échange
400\,000 actions par jour~?
\item Sa \omterm{prop:m1:what-a-trading-firm-does:capture}{captation nette} avant le snipe était de 0.35 centime par action. Est-il
toujours rentable~?
\item Il peut ramener son délai à 30 microsecondes pour 2.5 millions de dollars par an.
L'investissement en vaut-il la peine, sur 252 jours~?
\item Si les quatre \omterm{def:m1:what-a-trading-firm-does:market-maker}{teneurs de marché} réalisent cet investissement, que deviennent le profit
du trader et les coûts des \omterm{def:m1:what-a-trading-firm-does:market-maker}{teneurs de marché}~? Comment appelle-t-on
cette situation~?
\item Autrement, la plateforme impose un délai de 350 microsecondes aux ordres
entrants. Quelles cotations de ses \omterm{def:m1:what-a-trading-firm-does:market-maker}{teneurs de marché} sont-elles désormais sûres~?
\end{enumerate}

\medskip\noindent\textbf{Partie IV --- La règle.}
\begin{enumerate}[resume]
\item Un courtier achemine un ordre d'achat de son client en se fondant sur le flux. Pendant
un intervalle périmé après un saut à la hausse, où envoie-t-il l'ordre, et que
se passe-t-il~?
\item L'exécution constitue-t-elle un \omterm{def:m1:us-equity-market-structure:opr}{trade-through} au sens de la Rule 611~? Expliquer en
se référant à ce que chaque plateforme affiche à ce moment-là.
\item Les modifications de 2024 de la SEC réduiraient le \omterm{def:m1:us-equity-market-structure:cap}{plafond des frais d'accès} à 0.10 centime.
Refaire la question 9.
\item L'abrogation de la Rule 611 supprimerait-elle l'opportunité de snipe~?
\item Énoncer le \emph{résultat nommé}~: le profit espéré du trader par action
issu d'une cotation périmée, par saut, net des frais de preneur actuels, en centimes.
\item En une phrase, qui paie en fin de compte~?
\end{enumerate}
\end{problem}

\section{Questions d'entretien}

\begin{interviewq}[$\star$ \iqroles{trader, développeur, chercheur}]\label{iq:m1:us-equity-market-structure:1}
Qu'est-ce que le NBBO, qui le calcule, et pourquoi deux sociétés pourraient-elles ne pas s'accorder
sur sa valeur à une microseconde donnée~?
\end{interviewq}

\begin{interviewq}[$\star$ \iqroles{trader, développeur}]\label{iq:m1:us-equity-market-structure:2}
Qu'est-ce qu'un \omterm{def:m1:us-equity-market-structure:iso}{ordre de balayage intermarchés} (ISO), et pourquoi existe-t-il~?
\end{interviewq}

\begin{interviewq}[$\star\star$ \iqroles{trader, chercheur}]\label{iq:m1:us-equity-market-structure:3}
Pourquoi quelqu'un afficherait-il de la liquidité sur une plateforme qui fait payer
des frais aux makers, alors qu'une autre plateforme leur verse une rétrocession~?
\end{interviewq}

\begin{interviewq}[$\star\star$ \iqroles{développeur}]\label{iq:m1:us-equity-market-structure:4}
Vous devez maintenir le NBBO pour 8\,000 symboles à partir de seize \omterm{def:m1:us-equity-market-structure:sip}{flux directs}, avec
une latence minimale. Décrivez vos structures de données et votre modèle de threads.
\end{interviewq}

\begin{interviewq}[$\star\star$ \iqroles{chercheur, trader}]\label{iq:m1:us-equity-market-structure:5}
Plus de la moitié du volume des actions américaines s'échange hors bourse. Est-ce un problème~?
Pour qui~?
\end{interviewq}

\begin{interviewq}[$\star\star\star$ \iqroles{chercheur, trader}]\label{iq:m1:us-equity-market-structure:6}
Le régulateur propose d'abroger la \omterm{def:m1:us-equity-market-structure:opr}{règle de protection des ordres}. Prédisez trois
conséquences pour la structure de marché et dites comment vous positionneriez
une société de market making pour chacune d'entre elles.
\end{interviewq}
